% paper/sections/10_2_interface.tex
%
% §10.2  界面処理の精度検証
%   CLS 移流・再初期化精度，GFM 不連続面処理
% ═══════════════════════════════════════════════════════════════════════════════
\subsection{界面処理の精度検証}
\label{sec:verify_interface}
% ═══════════════════════════════════════════════════════════════════════════════

% ─────────────────────────────────────────────────────────────────────────────
\subsubsection{CLS 移流・再初期化精度}
\label{sec:verify_cls_advection}
% ─────────────────────────────────────────────────────────────────────────────

\noindent\textbf{目的：}
CLS 移流（§\ref{sec:advection}）と再初期化（§\ref{sec:reinit}）の
形状保持性能と質量保存性を標準テストで定量的に確認する．

\noindent\textbf{評価対象：}
DCCD 移流（$\varepsilon_d = 0.05$）＋再初期化（20 ステップ毎）の CLS パイプライン．

\noindent\textbf{手順とパラメタ：}
2 種類の標準テストを $N = 64, 128, 256$ で実施する．
\begin{itemize}[noitemsep]
  \item[(a)] \textbf{Zalesak スロット円板：}
        中心 $(0.5, 0.75)$，$R = 0.15$ の円板にスロット（幅 $0.05$，高さ $0.25$）を切り込み，
        中心 $(0.5, 0.5)$ の剛体回転場で 1 回転移流する．
  \item[(b)] \textbf{単一渦テスト（LeVeque 1996）：}
        円形界面（中心 $(0.5, 0.75)$，$R = 0.15$）を時間反転型渦場で
        $T/2$ まで強く変形させた後，$T/2 \to T$ で元に戻す．
\end{itemize}
評価指標：$L_2$/$L_\infty$ 形状誤差（$\psi - \psi_0$），相対質量保存誤差 $|\Delta M|/M_0$．

\noindent\textbf{結果：}

\begin{table}[htbp]
\centering
\caption{CLS 界面移流テストの精度（DCCD 移流 $\varepsilon_d = 0.05$）}
\label{tab:cls_advection}
\begin{tabular}{rrrrrr}
\toprule
 & $N$ & $L_2$ 誤差 & $L_\infty$ 誤差 & 質量保存誤差 \\
\midrule
\multirow{3}{*}{Zalesak}
 &  64 & $9.67 \times 10^{-2}$ & $5.06 \times 10^{-1}$ & $5.55 \times 10^{-3}$ \\
 & 128 & $9.28 \times 10^{-2}$ & $5.60 \times 10^{-1}$ & $1.08 \times 10^{-3}$ \\
 & 256 & $7.62 \times 10^{-2}$ & $5.63 \times 10^{-1}$ & $6.01 \times 10^{-5}$ \\
\midrule
\multirow{3}{*}{Single vortex}
 &  64 & $1.90 \times 10^{-1}$ & $8.46 \times 10^{-1}$ & $1.78 \times 10^{-3}$ \\
 & 128 & $1.73 \times 10^{-1}$ & $8.05 \times 10^{-1}$ & $1.63 \times 10^{-3}$ \\
 & 256 & $1.43 \times 10^{-1}$ & $7.44 \times 10^{-1}$ & $7.10 \times 10^{-5}$ \\
\bottomrule
\end{tabular}
\end{table}

$L_2$ 形状誤差は格子精緻化に伴い緩やかに減少した．
$L_\infty$ 誤差はスロット角部・界面不連続部で遷移幅 $\varepsilon = \Ord{h}$ が
精度の上界を規定するため単調減少しない場合がある．
質量保存誤差は $N = 256$ で $\Ord{10^{-5}}$ に達した．

\noindent\textbf{結論：}
DCCD 移流と再初期化の組み合わせは，質量保存誤差 $\Ord{10^{-5}}$（$N=256$）を達成し，
保存性が良好であることを確認した．
$L_\infty$ 形状誤差は界面厚み $\varepsilon$ により律速される．

% ─────────────────────────────────────────────────────────────────────────────
\subsubsection{GFM 不連続面処理の精度}
\label{sec:verify_gfm}
% ─────────────────────────────────────────────────────────────────────────────

\noindent\textbf{目的：}
GFM（§\ref{sec:gfm}）による圧力跳び処理が Young--Laplace 条件を
精度良く・粗格子でも安定に再現することを確認し，CSF との比較を行う．

\noindent\textbf{評価対象：}
GFM の PPE 右辺補正 $b^{\mathrm{GFM}}$（式~\eqref{eq:gfm_rhs_correction}）による
圧力跳び精度．CSF との比較．

\noindent\textbf{手順とパラメタ：}
\begin{itemize}[noitemsep]
  \item 計算領域 $[0,1]^2$，半径 $R = 0.25$ の円形液滴（$\rho_l/\rho_g = 1000$），静止状態
  \item 解析解 $\Delta p = \kappa/\mathit{We} = 1/(R \cdot \mathit{We})$
  \item 格子 $N = 32, 64, 128$
  \item $\varepsilon$ 感度評価：$\varepsilon/h = 0.5$--$4.0$（$N = 128$ で実施）
\end{itemize}

\noindent\textbf{結果：}

\begin{table}[htbp]
\centering
\caption{Laplace 圧力跳び $\Delta p = \kappa/\mathit{We}$（解析解 $4.0$）：GFM vs CSF}
\label{tab:gfm_laplace}
\begin{tabular}{rrrrr}
\toprule
$N$ & GFM $\Delta p$ & GFM 相対誤差 & CSF $\Delta p$ & CSF 相対誤差 \\
\midrule
 32 & 3.748 & $6.3 \times 10^{-2}$ & $-2.361$ & $1.6 \times 10^{0}$ \\
 64 & 3.866 & $3.4 \times 10^{-2}$ &  3.647  & $8.8 \times 10^{-2}$ \\
128 & 3.927 & $1.8 \times 10^{-2}$ &  3.952  & $1.2 \times 10^{-2}$ \\
\bottomrule
\end{tabular}
\end{table}

GFM は全格子で安定に収束し，$N = 128$ で相対誤差 $1.8\%$ を達成した（$\Ord{h^1}$）．
CSF は粗格子（$N = 32$）で $\Delta p$ の符号が反転したが，細格子では GFM を上回る精度に到達した．
$\varepsilon$ 感度評価では CCD 曲率が $\kappa \approx 4.007$（誤差 $< 0.2\%$）で
$\varepsilon$ 非依存を確認した．

\noindent\textbf{結論：}
GFM の収束次数は約 1 次（界面近傍の $\phi$ 補間精度が支配的）であり，
粗格子ロバスト性が大規模計算で有利である．

% ─────────────────────────────────────────────────────────────────────────────
\subsubsection{CLS 保存性評価：動的非一様格子上での保存型再マッピング}
\label{sec:cls_conservation}%
\label{app:cls_conservation}% backward compat
% ─────────────────────────────────────────────────────────────────────────────
% paper/sections/10_2a_cls_conservation.tex
%
% §10.2.3  CLS 法の保存性評価：動的非一様格子上での保存型再マッピング
%      実験スクリプト: experiments/ls_cls_conservation.py
%      結果データ:    results/ls_cls_conservation/

\noindent\textbf{目的：}
界面追従格子を定期的にリフレッシュする環境において，
CLS の保存型再マッピングが LS の非保存補間に比べ
質量誤差を低減することを定量的に検証する．

\noindent\textbf{評価対象：}
CLS 保存型再マッピング（質量スケーリング補正
$\psi_i \leftarrow \psi_i^{\mathrm{interp}} \cdot (M_{\mathrm{old}}/M_{\mathrm{new}})$）
vs LS 区分線形補間（非保存）．

\noindent\textbf{手順とパラメタ：}
\begin{itemize}[noitemsep]
  \item 格子 $N=128$，界面適合格子（圧縮比 $r=2$，$\sigma=0.06$），均一流 1 周期移流
  \item 空間微分 DCCD（$\varepsilon_d = 0.05$），時間積分 TVD-RK3，CFL $=0.4$
  \item $K_{\mathrm{refresh}} = 5, 10, 20, 50$ ステップごとに格子を界面位置に再生成
  \item 評価指標：$T=1$ 後の最終相対質量誤差 $|\Delta M/M_0|$
\end{itemize}

\noindent\textbf{結果：}

\begin{table}[htbp]
  \centering
  \caption{%
    $T=1$ 後の最終相対質量誤差 $|\Delta M/M_0|$（$N=128$，$r=2$，CFL $=0.4$）}
  \label{tab:ls_cls_mass_err}
  \renewcommand{\arraystretch}{1.3}
  \begin{tabular}{lrrrr}
    \toprule
    手法 & $K=5$ & $K=10$ & $K=20$ & $K=50$ \\
    \midrule
    LS (SDF)     & $8.6\times10^{-5}$ & $7.6\times10^{-5}$ & $1.2\times10^{-4}$ & $7.3\times10^{-4}$ \\
    CLS ($\psi$) & $7.7\times10^{-6}$ & $8.9\times10^{-7}$ & $7.4\times10^{-5}$ & $8.2\times10^{-4}$ \\
    \bottomrule
  \end{tabular}
\end{table}

$K=10$ で CLS は LS に比べ 85 倍の質量誤差低減を達成した
（CLS $8.9\times10^{-7}$ vs LS $7.6\times10^{-5}$）．
$K=5$ でも 11 倍の改善を示した．
大 $K$（$K=50$）では DCCD 移流誤差が支配的になり CLS の保存性優位が消えた．

\noindent\textbf{結論：}
CLS の保存型再マッピングはリフレッシュ起源の質量誤差を 1--2 桁低減する．
この優位性はリフレッシュ周期が短いほど顕著であり，
界面追従格子が頻繁に再生成される実用条件（$K \lesssim 20$）で特に有効である．

